\chapter{Stand der Technik}
\label{ch:standdertechnik}

\section{On-Device-Inferenz-Lösungen}
\label{sec:ondevice-loesungen}
Um ein Sprachmodell lokal auf einem Mobilgerät auszuführen, wird eine Laufzeitumgebung (Runtime) benötigt. Dafür stehen unterschiedliche Lösungen zur Verfügung. Eine solche Lösung ist die MediaPipe LLM Inference API \autocite{mediapipe_llm}. Für die technische Umsetzung werden sogenannte \enquote{Task Bundles} (Dateiformate wie .task oder .litertlm) genutzt \autocite{mediapipe_llm}. Diese Datenpakete bündeln das eigentliche Modell mit den notwendigen Tokenizer-Parametern (Werkzeuge zur Textzerlegung), um eine reibungslose Ende-zu-Ende-Verarbeitung zu gewährleisten \autocite{mediapipe_llm}.
\medskip

\noindent Abhängig vom Modelltyp und den Einstellungen können die Berechnungen entweder auf dem Hauptprozessor (CPU) oder dem Grafikprozessor (GPU) des mobilen Geräts durchgeführt werden \autocite{mediapipe_llm}.
\medskip

\noindent Die Schnittstelle unterstützt von Haus aus kompakte, speziell für mobile Geräte entwickelte Modelle. Dazu gehört insbesondere die Gemma-Familie von Google, etwa Gemma 3 1B und Gemma 2 2B \autocite{mediapipe_llm}. Für die folgenden Abschnitte sind vor allem diese beiden sehr kompakten Modelle wichtig.

\section{Die Gemma-Modellfamilie}
\label{sec:gemma}
Die Gemma-Modellfamilie gehört zu den offenen und ressourcenschonenden Modellen, die vom Gemma-Team bei Google DeepMind entwickelt wurden \autocite{gemma3_2025}.
\medskip
Besondere Merkmale der Gemma-Modelle sind:
\begin{itemize}
    \item \textbf{Offenheit und Effizienz:} Neben der freien Zugänglichkeit der Modelle zeichnen sie sich durch eine ressourcenschonende Ausführung auf herkömmlicher Standardhardware wie Smartphones aus \autocite{gemma3_2025,gemma2_2024}.
    \item \textbf{Skalierbarkeit:} Um unterschiedlichen technischen Anforderungen gerecht zu werden, sind die Modelle in mehreren Größen verfügbar. Diese werden in Milliarden von Parametern gemessen. So bietet die zweite Generation (Gemma 2) Varianten mit 2, 9 und 27 Milliarden Parametern \autocite{gemma2_2024}. Die dritte Generation (Gemma 3) erweitert diese Skalierung auf Modelle mit 1, 4, 12 und 27 Milliarden Parametern \autocite{gemma3_2025}.
    \item \textbf{Erweiterter Funktionsumfang:} Modelle vor der dritten Generation (Gemma 1--2) unterstützten nur reine Textverarbeitung. Seit der dritten Generation sind auch multimodale Eingaben möglich, sodass neben Text auch Bilder verarbeitet werden können \autocite{gemma3_2025}. Zudem wurden die Fähigkeiten zur Verarbeitung unterschiedlicher Sprachen sowie zum Erfassen deutlich längerer Textzusammenhänge verbessert \autocite{gemma3_2025}.
    \item \textbf{Einsatzbereitschaft:} Neben der allgemeinen Basisversion gibt es die Modelle auch in einer speziell angepassten Form, die gezielt für das Führen von Dialogen und das präzise Befolgen von Anweisungen trainiert wurde \autocite{gemma2_2024,gemma3_2025}.
\end{itemize}

\noindent In dieser Arbeit werden aus dieser Familie die kompakten Modelle Gemma 3 1B und Gemma 2 2B eingesetzt und verglichen.
% - Gemma = offene, kleine Google-Modelle
% - Gemma 2 \autocite{gemma2_2024}, Gemma 3 \autocite{gemma3_2025}
% - Warum: leichtgewichtig + Größen ermöglichen Vergleich 1B vs 2B
% - Gewählt: Gemma 3 1B + Gemma 2 2B

\section{Ausgangslage: bestehende Anwendung}
\label{sec:ausgangslage}
% - Unicorn Hybrid Black: 8 Kanäle, 250 Hz
% - Bestehende App: Aufnahme, Bluetooth -> Room-DB, Live-FFT
% - DEIN Beitrag: features/, llm/, analyse/ (Bestandscode nur 1 Button + 1 Zeile)